% =====================================================================
%  Lời mở đầu
% =====================================================================
\chapter*{Lời mở đầu}
\addcontentsline{toc}{chapter}{Lời mở đầu}
\markboth{Lời mở đầu}{Lời mở đầu}

Một hệ thống thông minh không bắt đầu từ thuật toán. Nó bắt đầu từ câu hỏi:
\emph{thế giới thực đang được viết lại thành những con số nào, và những con số
ấy có giữ lại đúng thứ cần giữ hay không}. Đề bài của Assignment 01 gọi tên
chuỗi ấy rất rõ --- Thế giới thực $\rightarrow$ Dữ liệu $\rightarrow$ Biểu diễn
$\rightarrow$ Mô hình $\rightarrow$ Dự đoán $\rightarrow$ Ứng dụng --- và yêu
cầu đi trọn chuỗi ấy chứ không dừng ở một mắt xích.

Báo cáo này trình bày hai hệ thống thông minh cỡ nhỏ, cùng đi trọn một chuỗi
như vậy trên hai miền khác hẳn nhau. Hệ thống thứ nhất là bài toán
\textbf{phân lớp} trong y tế: từ tám chỉ số sinh tồn của một người, ước lượng
nguy cơ mắc đái tháo đường. Hệ thống thứ hai là bài toán \textbf{hồi quy} trong
bất động sản: từ diện tích, kết cấu và vị trí hành chính của một căn nhà, ước
lượng giá giao dịch tính bằng tỷ đồng. Mỗi bài toán được giải bằng năm mô hình
Học máy truyền thống, so sánh dưới cùng một giao thức đánh giá, cộng thêm bốn
thực nghiệm có kiểm soát và một Đồ thị Tri thức năm tầng.

Điều mà Đồ thị Tri thức mang lại đáng được nói ngay từ đây, vì nó là phần mà
một mô hình học máy đơn lẻ không làm được. Mô hình trả về xác suất $54{,}5\%$;
đồ thị biến con số ấy thành một gói chăm sóc gồm năm mặt hàng có thật, chi phí
ban đầu $1{,}99$--$2{,}75$ triệu đồng và duy trì $1{,}20$--$1{,}66$ triệu đồng
mỗi tháng, mua được tại hệ thống nhà thuốc. Mô hình trả về giá $8{,}42$ tỷ
đồng; đồ thị biến con số ấy thành hạn mức vay $5{,}89$ tỷ, khoản trả góp
$47{,}5$ triệu đồng mỗi tháng trong 25 năm, và mức thu nhập tối thiểu
$118{,}7$ triệu đồng mỗi tháng để gánh được khoản vay ấy.

Việc đặt hai bài toán cạnh nhau không phải để nhân đôi khối lượng. Nó là một
thiết kế có chủ đích: khi cùng một thực nghiệm được lặp lại trên hai miền dữ
liệu, những kết luận nào còn đứng vững và những kết luận nào sụp đổ sẽ tự lộ
ra. Chương~\ref{ch:ketluan} cho thấy chính cách làm này đã \emph{bác bỏ} một
giả định phổ biến mà bản thân người viết cũng tin lúc mới bắt đầu --- rằng
``biểu diễn dữ liệu quan trọng hơn thuật toán''. Phát biểu ấy đúng ở bài toán
bất động sản và sai ở bài toán tiểu đường, và điều phân biệt hai trường hợp lại
không nằm ở thuật toán mà nằm ở chỗ phép biểu diễn có \emph{thêm thông tin mới}
hay chỉ \emph{đổi dạng} thông tin đã có.

Toàn bộ con số trong báo cáo được sinh ra từ một lệnh duy nhất
(\texttt{python run\_pipeline.py}) với hạt giống ngẫu nhiên cố định, nên mọi
bảng và mọi hình đều tái lập được. Mã nguồn, hai Jupyter Notebook, bộ 355 bài
kiểm thử và bản ứng dụng web đã triển khai công khai được liệt kê ở
Mục~\ref{sec:masource} ngay đầu Chương~\ref{ch:coso}.

Em xin chân thành cảm ơn thầy \textbf{PGS.TS Trần Đình Quế} đã hướng dẫn học
phần và đưa ra đề bài này --- một đề bài buộc người làm phải đi hết chuỗi xử lý
thay vì dừng lại ở chỗ mô hình cho ra con số đẹp.

\vspace{8mm}
\begin{flushright}
\emph{Hà Nội, tháng 8 năm 2026}\par
\vspace{2mm}
\textbf{Nguyễn Duy Nghĩa}
\end{flushright}
